\section{Recommended Final Demonstration}
\label{sec:final-demo}

A customer-facing November demonstration showing both adoption paths side by side, both built from hardware derived from the same Alfred common node (Section~\ref{sec:common-platform}).

\begin{diagram}
LEFT SIDE -- CLEAN-SHEET              RIGHT SIDE -- EXISTING VEHICLE

     Clicker 4                        legacy device -- CAN -- legacy device
         |                                          |
    10BASE-T1S                              Alfred Gateway
         |                                          |
  +------+------+                            10BASE-T1S
  |             |                                   |
 E2B           E2B                               Clicker 4
  |             |
sensor       actuator

           BOTH SIDES: SAME ALFRED COMMON NODE
\end{diagram}

The point of the demonstration is explicit and stated to the customer in exactly these terms: a customer does not have to choose whether Alfred only supports new machines or only supports legacy machines.

\begin{diagram}
NEW PROGRAM       -> native T1S + E2B
EXISTING PROGRAM  -> gateway + pseudo-bus

                  using the SAME underlying infrastructure
\end{diagram}

Run order: start with the right side (existing vehicle) so the customer sees their own architecture untouched and Alfred attaching non-invasively first -- this is deliberately the lower-anxiety opening for a brownfield-leaning audience. Show Phase~A (passive mirror) before Phase~B (authorized bidirectional command), narrating the Mode~1/Mode~2 distinction explicitly (Section~\ref{sec:gateway-modes}) so the safety posture is understood before any transmit capability is shown. Then move to the left side and show the identical Clicker~4 issuing a command to a native E2B actuator and receiving native E2B sensor data, closing on the single sentence that both sides just ran on hardware built from the same core design (Section~\ref{sec:common-platform}) -- that sentence is the actual sale, not the individual demos.
